Un canó d'electrons és un dispositiu format per un càtode i un ànode, i té per objectiu accelerar electrons que després seran dirigits mitjançant un camp magnètic per a bombardejar un blanc on es troba una mostra d'un cert material. La corba de la figura següent representa la trajectòria que segueix un electró des que surt del càtode fins que impacta en el blanc. L'objectiu és escalfar el material fins a poder-lo sublimar.

Els electrons són emesos pel càtode del canó i tenen una velocitat inicial nul·la. Posteriorment, els electrons són accelerats en l'espai de 20,0~cm que separa el càtode de l'ànode. El camp elèctric entre l'ànode i el càtode és uniforme i constant. La diferència de potencial entre el càtode i l'ànode és de 220~V.

\begin{apartat}{1,25}
Determineu el mòdul del camp elèctric que hi ha entre l'ànode i el càtode. Representeu sobre el dibuix les línies de camp elèctric indicant-ne clarament la direcció i el sentit. Quin punt està a un potencial més alt, l'ànode o el càtode? Justifiqueu les respostes.
\end{apartat}

\begin{apartat}{1,25}
Després de travessar l'ànode, els electrons són dirigits per un camp magnètic perquè impactin contra el blanc. En aquest espai, el camp elèctric és nul. Calculeu el mòdul de la velocitat dels electrons just després de travessar l'ànode i just abans d'impactar en el blanc. Quina direcció i quin sentit ha de tenir el camp magnètic perquè els electrons impactin en el blanc? Representeu sobre el dibuix el camp magnètic. Justifiqueu la resposta.
\end{apartat}

\dades{%
  $m_\mathrm{e} = 9,11\times 10^{-31}\ \mathrm{kg}$.\\
  $|e| = 1,602\times 10^{-19}\ \mathrm{C}$.}

\figura[0.5]{fig1.png}
